\section{What a mass breaks: duality, chirality, and the two senses of rotation}
\label{d2:sec:breaks}

The hypothesis reads $m$ as what breaks a symmetry of which $c^2$ is the symmetric part. The prior question is which symmetry a mass
actually breaks in the arena, and which it leaves alone. The arena has two things called mass. In the field-strength reading a frame vector
$X=\sum_i(B_ie_i+E_ie_{i+3})$, $z=\bB+i\bE$, has energy density $u=\half|z|^2$, complex invariant $q=z\cdot z$ and rest-frame energy density
$E_0=\half|q|$ (Section~14 of \cite{ES17}). In the Dirac reading the mass is the kernel $T=e_1e_2e_3=\gamma^0J$ of the Grassmann bilinear
(Proposition~4.5 and Corollary~18.4 of \cite{ES17}). In brief: a mass breaks no Lorentz generator and, in the Dirac reading, no generator of
$\so{3,3}$; it breaks the duality rotation generated by the binding root $G$, in the sense of \cref{d2:prop:duality,d2:thm:little}, and the
chiral and Pauli--G\"ursey symmetries of the massless kinetic term, down to $\uu1_J$; and the two senses of the Dirac rest rotation are exchanged by
charge conjugation, those of the frame's binding rotation by parity. Units are $\hbar=c=1$ unless a symbol is shown. Every finite-dimensional
statement was checked exactly by \texttt{verify/d2\_symmetry\_breaking.py}; the checks are listed in \cref{d2:tab:claims}.

\subsection{Duality and the rest energy of a field}
\label{d2:sec:duality}

Write $D_\vartheta=e^{\vartheta G}=\cos\vartheta\,\id+\sin\vartheta\,G$ (since $G^2=-\id$). As $G$ acts on $z$ as multiplication by $i$
(Theorem~12.1 of \cite{ES17}), $D_\vartheta$ is the duality rotation.

\begin{proposition}[Duality on the invariants]
\label{d2:prop:duality}
For every $X$ and $\vartheta$:
\begin{itemize}
\item[(i)] $z(D_\vartheta X)=e^{i\vartheta}z(X)$, $u(D_\vartheta X)=u(X)$, $q(D_\vartheta X)=e^{2i\vartheta}q(X)$;
\item[(ii)] $D_\vartheta^{\!\top}D_\vartheta=\id$ and $D_\vartheta^{\!\top}\eta\,D_\vartheta=\cos2\vartheta\,\eta+\sin2\vartheta\,\eta G$: $D_\vartheta$ preserves the positive form
for every $\vartheta$, and $\eta$ only for $\vartheta\in\pi\Z$;
\item[(iii)] if $q(X)=0$ every $\vartheta$ fixes $q$; if $q(X)\neq0$ exactly $\vartheta\in\pi\Z$ do;
\item[(iv)] if $q\neq0$, in a Lorentz frame where $z=(b+ia)\hat e_1$ with $a,b$ real, $q=(a^2+b^2)e^{2i\varphi}$ with $\tan\varphi=a/b$: $\arg q$ is
twice the Rainich complexion, and $D_{-\varphi}$ carries the field to the single magnetic plane $\sqrt{a^2+b^2}\,\theta^2\wedge\theta^3$.
\end{itemize}
\end{proposition}

\begin{proof}
(i) is $G\leftrightarrow i$ and the quadratic character of $q$. (ii) $G$ is antisymmetric, so $D_\vartheta$ is orthogonal; $G$ is
$\eta$-self-adjoint, so $G^{\!\top}=\eta G\eta$, and expanding $(\cos\vartheta+\sin\vartheta\,\eta G\eta)\eta(\cos\vartheta+\sin\vartheta\,G)$ with $G^2=-\id$
gives the formula; $\eta$ and $\eta G$ are linearly independent. (iii) is $e^{2i\vartheta}-1=2i\sin\vartheta\,e^{i\vartheta}$. (iv) The frame is that of
Proposition~12.12 of \cite{ES17}, $F=a\,\theta^0\wedge\theta^1+b\,\theta^2\wedge\theta^3$, so $z=(b+ia)\hat e_1$.
\end{proof}

Items (i), (iii) and (iv) are inherited, the classical duality rotation of the Maxwell invariants \cite{Rainich1925,Calkin65,DeserTeitelboim76}
written in the frame. Item (ii) is fixed by the arena: the generator of duality is the binding root, which lies in the $\eta$-self-adjoint summand of
$\End(V)=\so{3,3}\oplus\{X=X^\dagger\}$ (Theorem~17.3 of \cite{ES17}). For $\vartheta\notin\pi\Z$, $D_\vartheta\notin O(3,3)$, so it lifts to no element of
$\mathrm{Pin}(3,3)$ and has no action on spinors (as Section~17 of \cite{ES17} notes for the rotor). Duality is a symmetry of the positive form $u$; the
split-form interior Lagrangian is rotated by it into its $\eta G$ partner.

\begin{proposition}[The rest energy is the doubly invariant modulus]
\label{d2:prop:rest}
With $\mathbf S=\bE\times\bB$, $u^2-|\mathbf S|^2=\tfrac14|q|^2$, so $u\ge E_0$, with equality exactly when $\mathbf S=0$. Of $u$, $\Re q$, $\Im q$ and
$E_0$, the first is invariant under $U(1)_G$ but not under $\lG$, the next two under $\lG$ but not under $U(1)_G$, and $E_0=\half|q|$ under both.
\end{proposition}

\begin{proof}
$u^2-|\bE\times\bB|^2=\tfrac14(|\bB|^2+|\bE|^2)^2-|\bE|^2|\bB|^2+(\bE\cdot\bB)^2=\tfrac14(|\bB|^2-|\bE|^2)^2+(\bE\cdot\bB)^2$. The six generators of $\lG$
commute with $G$ and lie in $\so{3,3}$, so $\eta(X,AX)=\eta(X,GAX)=0$; a boost $K_i$ is symmetric for the positive form and $X^{\!\top}K_iX\neq0$ for
generic $X$. The duality statements are \cref{d2:prop:duality}.
\end{proof}

This is inherited; with units restored it reads $u^2-c^2|\mathbf g|^2=E_0^2$, $\mathbf g=\mathbf S/c^2$, the field form of $E^2-p^2c^2=m^2c^4$. Its
reading for the hypothesis is interpretive. Lorentz orbits of nonzero frame vectors are labelled by $q\in\C$, since $\lG\cong\so{3,\C}$ is
transitive on each level set of $q$ minus the origin (proof of Theorem~14.5 of \cite{ES17}); duality rotates the label through $2\vartheta$, fixes only
the origin, and the stabiliser of a nonzero label is $\{D_0,D_\pi\}=\{\pm\id\}$. A nonzero rest energy thus breaks $U(1)_G$ to $\Z_2$ on the space of
Lorentz orbits, and $\arg q$ is the angle along the broken direction. But the breaking is a property of a configuration, not of a vacuum ($u$ is
minimised by the zero field and no potential selects $|q|\neq0$, Section~14 of \cite{ES17}), and the broken symmetry is not one of the interior
Lagrangian. What the interior dynamics does with duality is sharper. Every solution of $\ddot X=-\kappa^2X$ (Theorem~16.2 of \cite{ES17}) is
$z(\tau)=c_-e^{-i\kappa\tau}+c_+e^{i\kappa\tau}$, $c_\pm\in\C^3$, and the zero-momentum orbits are those with $c_+=0$ (Proposition~16.5 of \cite{ES17}).

\begin{proposition}[The interior energy carries duality charge two]
\label{d2:prop:interior}
Along every solution:
\begin{itemize}
\item[(i)] $\mathcal E_\C:=\tfrac m2\bigl(q(\dot z)+\kappa^2q(z)\bigr)=2m\kappa^2\,c_-\cdot c_+$ is conserved, and the Noether energy
$\tfrac m2\eta(\dot X,\dot X)-\tfrac{m\kappa^2}2\eta(GX,GX)$ is $\Re\mathcal E_\C$; it vanishes when either sense is absent;
\item[(ii)] a constant duality rotation maps solutions to solutions and multiplies $\mathcal E_\C$ by $e^{2i\vartheta}$;
\item[(iii)] $p=m\eta(\dot X+\kappa GX)$ has $|p|^2=4m^2\kappa^2|c_+|^2$, while $u(\tau)=\half(|c_-|^2+|c_+|^2)+\Re\bigl(\bar c_-\cdot c_+\,e^{2i\kappa\tau}\bigr)$
is constant exactly when $\bar c_-\cdot c_+=0$;
\item[(iv)] on a zero-momentum orbit $p=0$ and the Noether energy is $0$, while $u=\half|c_-|^2$ and $E_0=\half|c_-\cdot c_-|$ are constant: the motion
$z=e^{-i\kappa\tau}c_-$ is a duality rotation along $\arg q$.
\end{itemize}
\end{proposition}

\begin{proof}
$\dot z+i\kappa z=2i\kappa c_+e^{i\kappa\tau}$ is the image of $\alpha_+X$, which gives $|p|^2$; in $q(\dot z)+\kappa^2q(z)$ the terms in $e^{\mp2i\kappa\tau}$
cancel, leaving $4\kappa^2c_-\cdot c_+$; $\eta(GX,GX)=-\eta(X,X)$ identifies the Noether energy with $\tfrac m2\Re(q(\dot z)+\kappa^2q(z))$; $q$ has duality
charge two; and on $c_+=0$, $\dot X=-\kappa GX$. Each identity was also checked exactly as a polynomial identity in the real and imaginary parts of
$c_\pm$, modulo $\cos^2\kappa\tau+\sin^2\kappa\tau=1$.
\end{proof}

The proposition is fixed; (i) restates Proposition~16.5(iii) of \cite{ES17}, and (ii)--(iv) are added. It corrects the working dictionary. The
duality-invariant $u=\half|z|^2$ is not a constant of the interior motion, and it is not equivalent to the closure Casimir $|p|^2$: on the
zero-momentum orbits, exactly those along which duality acts, $|p|^2=0$ while $u>0$. The two agree only if the field is read off the momentum, as in
the closure variable $w(p)$ of Section~17 of \cite{ES17}. And the interior's conserved energy is not undivided: it is the real part of a constant of
duality charge two, carried by neither sense alone.

\subsection{Little groups: field points and momenta}
\label{d2:sec:little}

On $z$, $\omega\cdot L+\mathbf v\cdot K\in\lG$ acts as $z\mapsto(\omega+i\mathbf v)\times z$, that is $\delta\bE=\mathbf v\times\bB$, $\delta\bB=-\mathbf v\times\bE$
(Theorem~12.3 of \cite{ES17}): $\lG=\so{3,\C}$ acting by complex cross products, and $G$ acts as $i$.

\begin{theorem}[Stabilisers of field points]
\label{d2:thm:little}
Let $z\neq0$.
\begin{itemize}
\item[(i)] $\Stab_{\lG}(z)=\{(\lambda z)\times\cdot:\lambda\in\C\}$ is abelian of real dimension two. If $q\neq0$ and $a=z/\sqrt q$, then $R=a\times\cdot$
has $R^3=-R$ (a rotation) and $(GR)^3=GR$ (a boost) about the common complex axis $a$: $\so2\oplus\so{1,1}$. If $q=0$, every element $N$ of the
stabiliser has $N^3=0\neq N^2$: null rotations, $\R^2$.
\item[(ii)] In $\lG\oplus\R G$ a non-null $z$ has the same stabiliser, with no $G$ component. A null $z=\mathbf u+i\mathbf v$ has a three-dimensional
stabiliser $\cong\mathfrak e(2)$: the null rotations and $L_{\mathbf n}+G$, with $\mathbf n=\mathbf u\times\mathbf v/|\mathbf u|^2$ the propagation axis,
acting on the null rotations with eigenvalues $\pm i$. For $z=(1,i,0)$, $(L_3+G)z=0$ and $e^{\vartheta G}z=e^{-\vartheta L_3}z$.
\item[(iii)] In $\lint$ on $E=\langle e_1,e_4,e_5,e_6\rangle$, $\Stab(e_1)=\so3$ and $\Stab(e_1+e_6)=\mathfrak e(2)$, both three-dimensional.
\end{itemize}
\end{theorem}

\begin{proof}
(i) $a\times z=0$ iff $a\in\C z$. From $(z\times)^2w=z(z\cdot w)-qw$, $(z\times)^3=-q\,(z\times)$; with $a\cdot a=1$ this is $R^3=-R$, and $G\leftrightarrow i$
gives $(GR)^3=GR$; for $q=0$, $(z\times)^3=0\neq(z\times)^2$. (ii) The condition is $a\times z+i\vartheta z=0$; dotting with $z$ gives $i\vartheta q=0$.
For null $z$, $(\mathbf u\times\mathbf v)\times\mathbf u=|\mathbf u|^2\mathbf v$ and $(\mathbf u\times\mathbf v)\times\mathbf v=-|\mathbf u|^2\mathbf u$ give
$\mathbf n\times z=-iz$, so $(L_{\mathbf n}+G)z=0$, and $[L_{\mathbf n},z\times\cdot]=(\mathbf n\times z)\times\cdot=-i(z\times\cdot)$. (iii) is the classical
computation. All items were also verified exactly at $z=(1+i,1,i)$, $q=2i$, and at the null vectors $(1,i,0)$ and $(1+2i,2+i,2-2i)$.
\end{proof}

The theorem is inherited: (iii) is Wigner's classification \cite{Wigner39,WeinbergQFT}, (i)--(ii) the classical stabilisers of a bivector and the
identity of helicity and duality on null fields \cite{PenroseRindler84,Calkin65}, recorded on the frame in Example~13.11 of \cite{ES17}. It separates
two notions the hypothesis runs together. A particle's little group stabilises its momentum and is three-dimensional with or without mass; a field
point's stabiliser is two-dimensional with or without $q$. A mass changes the type of a Lorentz stabiliser, never its dimension, and breaks no Lorentz
generator of the theory, since $q$ and $P^2$ are invariants. What a nonzero $q$ removes from the stabiliser is the duality generator: a null field
point is duality-invariant up to a rotation about its axis, so Wigner's $\mathfrak e(2)$ reappears once duality is admitted, with helicity as compact
generator, while for a non-null point duality is broken to $\{\pm\id\}$.

\subsection{The Dirac mass: invariant in six dimensions, breaking the kinetic algebra}
\label{d2:sec:dirac}

With $\gamma^0=e_1$, $\gamma^a=e_{a+3}$ the kinetic kernels are $K^\mu=-\gamma^0\gamma^\mu$ and the mass kernel is $T=\gamma^0J$ (Proposition~4.5 of
\cite{ES17}); a Grassmann bilinear $\psi^{\!\top}M\psi$ is invariant under $\delta\psi=X\psi$ iff $X^{\!\top}M+MX=0$.

\begin{proposition}[The mass form is a $\mathrm{Spin}(3,3)$ invariant; the kinetic form is not]
\label{d2:prop:Tinv}
\begin{itemize}
\item[(i)] Every one of the fifteen bivectors $X=e_ae_b$ satisfies $X^{\!\top}T+TX=0$, and every one-parameter subgroup they generate
satisfies $S^{\!\top}TS=T$. The antisymmetric $M\in M_8(\R)$ satisfying the fifteen conditions form exactly $\R\,T$.
\item[(ii)] An element $X\in\so{3,3}$ satisfies $X^{\!\top}K^\mu+K^\mu X=\lambda^\mu{}_\nu K^\nu$ for some real $4\times4$ matrix $\lambda$ iff
$X\in\lint\oplus\R J$, of dimension $7$.
\end{itemize}
\end{proposition}

\begin{proof}
(i) By Lemma~2.1 of \cite{ES17} the trivector $T$ commutes with a bivector sharing two or no indices with $\{1,2,3\}$ and anticommutes with a mixed
one; by Theorem~2.3 the former are the antisymmetric and the latter the symmetric bivectors, so $X^{\!\top}T+TX$ is $[T,X]$ or $\{T,X\}$ and vanishes.
On the subgroups, $\cos\theta+\sin\theta X$ commutes with $T$, and $S=\cosh\theta+\sinh\theta X$ gives $S^{\!\top}TS=TS^{-1}S=T$. Uniqueness is
Corollary~18.4 of \cite{ES17}, re-derived by solving the $960$ linear conditions on the $28$ antisymmetric entries. (ii) The $256$ linear equations in
the fifteen bivector coefficients and the sixteen entries of $\lambda$ have a seven-dimensional solution space, projecting injectively to the
bivectors; $\lint$ and $J$ lie in it by covariance and by $[J,\gamma^0\gamma^\mu]=0$.
\end{proof}

Both statements are fixed. Within $\mathrm{Spin}(3,3)$ the Dirac mass breaks nothing, and what reduces the six-dimensional symmetry to the interior
one is the kinetic term, where the clock is converted into the spatial triple and $c$ enters. In the six-dimensional sense the entry \emph{$c^2$
symmetric, $m$ breaking} is inverted; the breaking it describes lives in the internal symmetries of the kinetic term.

\begin{proposition}[What the Dirac mass breaks]
\label{d2:prop:massless}
\begin{itemize}
\item[(i)] The solutions $X\in\End_\R(\R^8)$ of $X^{\!\top}K^\mu+K^\mu X=0$ ($\mu=0,\dots,3$) form
$\mathfrak g_{\rm massless}=\operatorname{span}\{\Omega,J,e_2\Omega,e_3\Omega\}=\uu1_\Omega\oplus\su2$, with trace form $-8\,\id$; adding
$X^{\!\top}T+TX=0$ leaves $\R J$.
\item[(ii)] On $\mathfrak g_{\rm massless}$ the orbit map $X\mapsto X^{\!\top}T+TX$ has rank $3$ and image $\operatorname{span}\{T\Omega,Te_2\Omega,Te_3\Omega\}$,
inside the six-dimensional space of $\lint$-invariant kernels; $T\Omega$ commutes with $J$ (pseudoscalar mass), $Te_2\Omega$ and $Te_3\Omega$ anticommute
with it (Majorana kernels).
\item[(iii)] For $M=T(m_s+m_p\Omega)$ the field equation $K^\mu\partial_\mu\Psi=M\Psi$ of $\half\Psi^{\!\top}(K^\mu\partial_\mu-M)\Psi$, with
$\partial_a\Psi=Jp_a\Psi$, is $\dot\Psi=-JH\Psi$ where
\begin{equation}
\label{d2:eq:Hmp}
H=\gamma^0\gamma^ap_a+\gamma^0(m_s+m_p\Omega),\qquad H^{\!\top}=H,\qquad [H,J]=0,\qquad H^2=\bigl(|\bp|^2+m_s^2+m_p^2\bigr)\id .
\end{equation}
\item[(iv)] $e^{\alpha\Omega}$ preserves each $K^\mu$ and sends $M\mapsto Me^{2\alpha\Omega}$; the stabiliser of $M\neq0$ in $U(1)_\Omega$ is $\{\pm\id\}$.
\end{itemize}
\end{proposition}

\begin{proof}
(i) The $256$ conditions on the $64$ entries of $X$ have exactly these solutions (Proposition~5.5 of \cite{ES17}, recomputed). At $\mu=0$ the condition
is $X^{\!\top}=-X$, so mass invariance is $[T,X]=0$, which only $J$ satisfies. (ii) For antisymmetric $X$ the map is $[T,X]$; by the blade rule $T$
anticommutes with $\Omega$, $e_2\Omega$, $e_3\Omega$ and commutes with $J$, giving $2T\Omega,2Te_2\Omega,2Te_3\Omega$; $e_2,e_3$ anticommute with $J$, $T$
and $\Omega$ commute with it. (iii) Varying gives $K^\mu\partial_\mu\Psi=M\Psi$, since $K^\mu$ is symmetric and $M$ antisymmetric. With $K^0=-\id$ and $J$
commuting with $\gamma^0\gamma^a$, $\gamma^0$, $\Omega$: $\dot\Psi=(K^aJp_a-M)\Psi=-J(\gamma^0\gamma^ap_a+\gamma^0(m_s+m_p\Omega))\Psi$. The five matrices
$\gamma^0\gamma^a$, $\gamma^0$, $\gamma^0\Omega=e_4e_5e_6$ are symmetric, square to $\id$ and pairwise anticommute, which gives $H^2$. (iv) $\Omega\in\mathfrak
g_{\rm massless}$ and $\{T,\Omega\}=0$, so $e^{-\alpha\Omega}Te^{\alpha\Omega}=Te^{2\alpha\Omega}$.
\end{proof}

Items (i) and (ii) are fixed, re-verifying Proposition~5.5 and Theorem~18.5(a) of \cite{ES17} and naming the three broken directions; their physical
content, that a Dirac mass breaks the axial $U(1)$ and the Pauli--G\"ursey $SU(2)$ of a massless Dirac field to fermion number \cite{Pauli57,Gursey58},
is inherited by Theorem~4.3 of \cite{ES17}. Items (iii) and (iv) are inherited, the realification of $h=\boldsymbol\alpha\cdot\bp+\beta(m_s+i\gamma_5m_p)$
and of the chiral rotation; (iii) was derived from the Grassmann kernels, and the commuting first-order form
$\half\Psi^{\!\top}J\dot\Psi-\half\Psi^{\!\top}H\Psi$ of Section~4 of \cite{ES17} with the same $H$ has the same flow. The parallel with the frame is exact at the level of identities: in both readings the mass is the modulus of an order
parameter of charge two under a $U(1)$ generated by a complex structure commuting with the Lorentz algebra ($G$ for $\lG$, $\Omega$ for $\lint$), and a
nonzero modulus leaves $\{\pm\id\}$. It stops there: the Dirac mass is a coupling, or a value selected by a potential (Theorem~18.5 of \cite{ES17}); the
rest energy of a frame vector is a property of a configuration that no potential selects.

\subsection{Particle and antiparticle: the two senses of one rotation}
\label{d2:sec:senses}

At $\bp=0$ the real first-order form $\half\Psi^{\!\top}J\dot\Psi-\half\Psi^{\!\top}H\Psi$ of Section~4 of \cite{ES17} has $H(0)=m\gamma^0=me_1$.

\begin{proposition}[The rest flow]
\label{d2:prop:restflow}
$JH(0)=mT$, so $\Psi(\lambda)=e^{-m\lambda T}\Psi_0=(\cos m\lambda-\sin m\lambda\,T)\Psi_0$. On the eigenspaces $e_1=\pm1$, each four-dimensional and
$J$-invariant, $T=\pm J$: $\Psi(\lambda)=e^{\mp m\lambda J}\Psi_0$, with energy $\half\Psi^{\!\top}H\Psi=\pm\half m|\Psi|^2$. Moreover $\gamma^0=-TJ$.
\end{proposition}

\begin{proof}
$Je_1=e_1e_2e_3$; $T=e_1J$ with $[e_1,J]=0$ gives $TP_\pm=\pm JP_\pm$ for $P_\pm=\half(\id\pm e_1)$; $T^2=-\id$ gives the exponential; $-TJ=-e_1J^2=e_1$.
\end{proof}

This is inherited, the realification of $\psi=e^{\mp imt}\psi_0$ on $\beta=\pm1$, at angular frequency $mc^2/\hbar$ in restored units. The arena adds a
name: the rest motion is the one-parameter group of the temporal volume element $T$, the unique $\mathrm{Spin}(3,3)$-invariant mass form, and the clock
$\gamma^0=-TJ$ is the quotient of that rotation by the charge rotation. Read interpretively, the single rotation $e^{-m\lambda T}$ splits into two
senses of the charge rotation according to the eigenvalue of the clock.

\begin{proposition}[Which operation reverses the sense]
\label{d2:prop:senses}
\begin{itemize}
\item[(i)] Every unit vector $C$ of the surplus plane $U$ has $C^2=\id$, $[C,T]=0$, $\{C,J\}=\{C,e_1\}=0$ and $C^{\!\top}H(0)C=-H(0)$; it maps the solution
$e^{-m\lambda J}\Psi_0$ on $e_1=+1$ to the solution $e^{+m\lambda J}C\Psi_0$ on $e_1=-1$ of the same flow. The parities $\gamma^0$, $\gamma^0J$ and the
time reversal $T_{\rm rev}=e_2e_4e_5e_6$ (Theorem~5.3 of \cite{ES17}) keep the sign of the rest energy.
\item[(ii)] On the frame, parity ($\bE\mapsto-\bE$) acts as $P=\eta$, so $PGP^{-1}=-G$, $\mathcal L_\kappa(PX)=\mathcal L_{-\kappa}(X)$ and
$\mathcal L_\kappa-\mathcal L_{-\kappa}=-2m\kappa\,\tfrac{d}{d\tau}(\bE\cdot\bB)$, while $X\mapsto-X$ and $X(\tau)\mapsto-\eta X(-\tau)$ leave
$\mathcal L_\kappa=\tfrac m2\eta(\alpha_+X,\alpha_+X)$ invariant. On $z=c_-e^{-i\kappa\tau}+c_+e^{i\kappa\tau}$:
$P\colon(c_-,c_+)\mapsto(\bar c_+,\bar c_-)$, $C\colon(c_-,c_+)\mapsto(-c_-,-c_+)$, $T\colon(c_-,c_+)\mapsto(-\bar c_-,-\bar c_+)$.
\item[(iii)] The sense of the Dirac rest rotation is reversed by $C$, $CP$, $CT$, $CPT$; that of the frame's binding rotation by $P$, $CP$, $PT$, $CPT$.
\end{itemize}
\end{proposition}

\begin{proof}
(i) $C=\cos\phi\,e_2+\sin\phi\,e_3$; by the blade rule it anticommutes with $e_1$ and $J$ and commutes with $T$, so $C^{\!\top}e_1C=-e_1$ and $C$ commutes
with $e^{-m\lambda T}$. $\gamma^0$, $\gamma^0J$ commute with $e_1$; $T_{\rm rev}$ commutes with $e_1$ and anticommutes with $J$ and $T$, so
$T_{\rm rev}\Psi(-\lambda)$ solves the same flow on the same eigenspace. (ii) $\eta G\eta=G^{\!\top}=-G$ and $\eta$ is an $\eta$-isometry, so
$\mathcal L_\kappa(PX)=\tfrac m2\eta(\dot X-\kappa GX,\dot X-\kappa GX)$; the total derivative is Theorem~16.6 of \cite{ES17}; for $T$ the reversal of $\tau$
compensates the sign. On $z$: $P\colon z\mapsto\bar z$, $C\colon z\mapsto-z$, $T\colon z(\tau)\mapsto-\overline{z(-\tau)}$. (iii) Signs multiply under
composition (also checked on the matrices).
\end{proof}

Item (i) is inherited, realised by the arena's $C$ (Theorem~5.1 of \cite{ES17}); (ii) and (iii) are fixed, (ii) re-verifying and extending
Theorem~16.6 of \cite{ES17}. For Dirac matter the two senses are particle and antiparticle: solutions of one Lagrangian, of angular frequency $m$, each
with energy of definite sign, exchanged by the symmetry $C$. For the frame the two senses are the kernels of the two factors of
$\alpha_-\alpha_+=\partial_\tau^2+\kappa^2$: one carries no momentum, the other all of it; each alone carries zero energy (\cref{d2:prop:interior}); their
frequency $\kappa=c/\ell$ is a conversion rate, not a mass; and they are exchanged by parity, a symmetry only up to a total derivative. That the
frame's binding is not a particle--antiparticle pair in the Dirac sense is an interpretive conclusion from these fixed facts. In both readings the
reversal is $CPT$-odd.

\subsection{Two vacua}
\label{d2:sec:vacua}

\begin{proposition}[The Fock vacuum of one momentum]
\label{d2:prop:fock}
Let $\theta_1,\dots,\theta_8$ be antisymmetric $16\times16$ integer matrices with $\theta_i^2=-\id$, pairwise anticommuting, and
$\rho(A)=-\tfrac14\sum A_{ij}\theta_i\theta_j$. Then $\rho$ is a Lie homomorphism with $[\rho(A),\theta_k]=\sum_jA_{jk}\theta_j$;
$\chi_{2\rho(J)}=x^6(x^2+4)^4(x^2+16)$, $\chi_{2\rho(JH)}=x^6(x^2+4E^2)^4(x^2+16E^2)$; the joint spectrum of $Q=-i\rho(J)$ and $H=-i\rho(JH(\bp))$ is the
table of Theorem~5.7 of \cite{ES17}, with a unique charge-zero vacuum at $-2E$ and multiplicities $1,4,6,4,1$ for $H+2E$. At $\bp=0$ the vacuum energy is
$-2m$ ($-2mc^2$ per momentum mode), minus half the four rest energies. The charge conjugation $C=e_2$ lifts to $U=\theta(v_1)\cdots\theta(v_4)/4$, with
$v_k$ orthogonal integer vectors, $|v_k|^2=2$, spanning the $-1$ eigenspace of $C$; $U$ is orthogonal, $U^2=\id$, $UQU^{-1}=-Q$, $UH(\bp)U^{-1}=H(-\bp)$.
\end{proposition}

\begin{proof}
The $\theta_i$ are tensor words in $\sigma_x,\sigma_z,\varepsilon$ found by exhaustive search; their $256$ monomials have rank $256$, so their commutant
is $\R\id$ (Theorem~5.6 of \cite{ES17}). The characteristic polynomials were computed exactly at $(\bp,m,E)=(0,1,1)$, $((0,0,3),4,5)$, $((1,2,2),0,3)$,
$((2,3,6),0,7)$, and, at the first two, the joint eigenspaces as kernels of the real systems equivalent to $(Q-q)v=(H-h)v=0$. A reflection in $v$ lifts to
$\theta(v)$ up to sign, and four signs cancel; $U\rho(A)U^{-1}=\rho(CAC^{\!\top})$, with $CJC^{\!\top}=-J$ and $CJH(\bp)C^{\!\top}=JH(-\bp)$ (Theorem~5.1
of \cite{ES17}).
\end{proof}

This re-verifies Theorem~5.7 and Proposition~5.8 of \cite{ES17} and is inherited: the Fock space of four fermionic modes of one Dirac momentum. The
phrase \emph{the vacuum carries the rest energy with the opposite sign} is accurate for it, and only for it.

\begin{proposition}[The frame vacuum]
\label{d2:prop:framevac}
The interior energy is $\mathcal E=\tfrac m2\bigl(\eta(\dot X,\dot X)+\kappa^2\eta(X,X)\bigr)=\sum_a\eta_{aa}\,\tfrac m2(\dot x_a^2+\kappa^2x_a^2)$; a spatial
mode has $p_y=-m\dot y$ and Hamiltonian $-(p_y^2/2m+m\kappa^2y^2/2)$. With canonical, positive-norm quantisation a temporal mode contributes
$\hbar\kappa(n+\half)$ and a spatial one $-\hbar\kappa(n+\half)$, so on a coordinate subspace of signature $(p,q)$ the spectrum is
$\hbar\kappa(\sum_Tn_a-\sum_Sn_a)+\tfrac{\hbar\kappa}2(p-q)$: the zero-point energy is $0$ on $\R^{3,3}$, $\pm\tfrac32\hbar\kappa$ on the temporal and
spatial triples, $\pm\half\hbar\kappa$ on the $(2,1)$ and $(1,2)$ planes.
\end{proposition}

\begin{proof}
The first identity uses $\eta(GX,GX)=-\eta(X,X)$; the spatial Hamiltonian is the Legendre transform of $-\tfrac m2(\dot y^2-\kappa^2y^2)$. With
$x=\sqrt{\hbar/2m\kappa}\,(a+a^\dagger)$, $p=i\sqrt{\hbar m\kappa/2}\,(a^\dagger-a)$, $p^2/2m+m\kappa^2x^2/2=\tfrac{\hbar\kappa}2(aa^\dagger+a^\dagger a)$,
checked exactly on truncated ladder matrices. The cross term is a total derivative; it shifts the momenta by a canonical transformation and leaves the
spectrum unchanged.
\end{proof}

\begin{corollary}[No ground state]
\label{d2:cor:noground}
With positive norm the spectrum of the full frame is $\hbar\kappa\Z$, every level is infinitely degenerate, and every state has one of lower energy
(add a spatial quantum). Making the spatial energies positive instead requires negative norm (Section~18 of \cite{ES17}).
\end{corollary}

The proposition is fixed, re-verifying Theorem~18.6 of \cite{ES17}: the zero-point energy vanishes because $P=Q$, the condition under which
factorisation and binding coexist (Theorem~16.2 of \cite{ES17}). The corollary is an obstruction under the hypothesis of positive norm, the familiar
fate of indefinite oscillators \cite{PaisUhlenbeck50}. The two vacua differ in kind: the Dirac vacuum has energy $-2E$ per momentum and is bounded
below; the frame's zero-point energy cancels between positive temporal and negative spatial oscillators, and there is no ground state for it to be
the energy of.

\subsection{Goldstone bookkeeping across the tower}
\label{d2:sec:goldstone}

\begin{table}[t]
\centering
\small
\begin{tabular}{@{}lllll@{}}
\toprule
stage & order parameter & kept & broken directions & status here\\
\midrule
$n=1$, Dirac & $T=\gamma^0J$ & $\uu1_J$ & $3$: $T\Omega$, $Te_2\Omega$, $Te_3\Omega$ (global) & re-verified\\
$n=1$, Majorana & $M=\Gamma M=M\Gamma$ & $\R(J-J\Gamma)$ & $3$ (global) & re-verified\\
$n=3$, Pati--Salam & $\Phi_0$ ($\nu_R\nu_R$) & $14$ of $23$ & $9$, eaten by $Z'$, $W_R^\pm$, six $X$ & re-verified\\
$n=3$, electroweak & $\Phi_0$ + bidoublet & $10$ of $23$ & $12$ eaten, $1$ axion-like ($1-\Gamma$) & count cited\\
\bottomrule
\end{tabular}
\caption{Broken directions attached to the specified generators: at $n=1$ the four of $\mathfrak g_{\rm massless}$, at $n=3$ the $21$
Pati--Salam generators and the central $J$, $J\Gamma$ (Proposition~22.1 of \cite{ES17}).}
\label{d2:tab:goldstone}
\end{table}

\begin{proposition}[Goldstone directions]
\label{d2:prop:goldstone}
\begin{itemize}
\item[(i)] The $\lint$-invariant kernels supported on $\Gamma=+1$ form a complex line; each $M$ in it anticommutes with $J$, has rank four,
$M^{\!\top}M\propto\half(\id+\Gamma)$ and $\Tr(M^{\!\top}M)^2/(\Tr M^{\!\top}M)^2=\tfrac14$; its stabiliser in $\mathfrak g_{\rm massless}$ is
$\R(J-J\Gamma)=\R(J+\Omega)$ and its orbit is three-dimensional.
\item[(ii)] At three constituents, with $P_*$, $\Phi_0$, $t^R_3$ of Appendix~A of \cite{ES17}: $[J,\Phi_0]=[J\Gamma,\Phi_0]=2[t^R_3,\Phi_0]\neq0$, and the
orbits of the $21$ Pati--Salam generators and of all $23$ have tangent spaces of the same dimension $9$ at $\Phi_0$.
\item[(iii)] If a kernel $K$ commutes with $J$ and anticommutes with $\Gamma$ (Dirac type), $[J-J\Gamma,K]=-2J\Gamma K\neq0$; if it lives on $\Gamma=+1$ and
anticommutes with $J$, $[J-J\Gamma,K]=0$. At one constituent $J-J\Gamma=J+\Omega$.
\end{itemize}
\end{proposition}

\begin{proof}
(i) and (ii) are exact computations, at three constituents on $\R^{512}$ with dyadic matrices and ranks over $\Q$; the reconstruction reproduces
Theorem~19.2(v) of \cite{ES17}: $\Phi_0$ antisymmetric, $\{J,\Phi_0\}=0$, $\Phi_0^{\!\top}\Phi_0$ the rank-four $\nu_R$ projector of the block,
$[t^R_3,\Phi_0]=J\Phi_0$. The identity in (ii) has a one-line reason: $\Phi_0$ lives on $\Gamma=+1$, so $J\Gamma\Phi_0=J\Phi_0$ and
$\Phi_0J\Gamma=-J\Phi_0$. (iii) $[J,\Gamma]=0$; if $\{K,\Gamma\}=0=[K,J]$ then $KJ\Gamma=-J\Gamma K$; if $K$ lives on $\Gamma=+1$ and $\{K,J\}=0$ then
$[J,K]=2JK=[J\Gamma,K]$; and $J\Gamma=J^2\Omega=-\Omega$. The bidoublet kernel $P_*E_1E_7E_8P_*$ served as an explicit Dirac-type instance.
\end{proof}

Items (i) and (ii) are fixed, re-verifying Theorems~18.5(b) and~22.2 of \cite{ES17}: the would-be Majoron is eaten. Item (iii) is fixed and is the
cross-scale statement: one element, $J-J\Gamma$, of charge $1-\Gamma$, is kept by a Majorana mass on $\Gamma=+1$ and broken by every Dirac-type mass, at one
constituent, where it is $J+\Omega$, and at three. That at the generic electroweak vacuum the broken $1-\Gamma$ direction escapes the gauge orbit
(Theorem~22.3 of \cite{ES17}) is cited, not re-verified.

\subsection{What the hypothesis gets right, and where the $b$-field belongs}
\label{d2:sec:verdict}

\begin{table}[t]
\centering
\small
\begin{tabular}{@{}llll@{}}
\toprule
symmetry & generator & broken by a mass? & class\\
\midrule
Lorentz $\lint$ & bivectors of $E$ & no: $P^2$, $T$ invariant; little groups stay $3$-dim & inherited\\
Lorentz $\lG$ & $L_i$, $K_i$ & no: $q$ invariant; stabilisers stay $2$-dim & inherited\\
$\mathrm{Spin}(3,3)$ & $15$ bivectors & no: $T$ invariant (the kinetic term keeps $7$) & fixed\\
duality $U(1)_G$ & $G$ & yes, to $\{\pm\id\}$ on Lorentz orbits, iff $q\neq0$ & inherited\\
axial $U(1)_\Omega$ & $\Omega$ & yes, to $\{\pm\id\}$ & inherited\\
Pauli--G\"ursey $\su2$ & $J$, $e_2\Omega$, $e_3\Omega$ & yes, to $\uu1_J$ & fixed\\
$U(1)_{1-\Gamma}$ & $J-J\Gamma$ & by Dirac type; not by Majorana on $\Gamma=+1$ & fixed\\
\bottomrule
\end{tabular}
\caption{What a mass breaks in the arena.}
\label{d2:tab:verdict}
\end{table}

\Cref{d2:tab:verdict} collects the answer. \emph{$m$ breaks the symmetry} is true for duality in the frame and for the chiral and Pauli--G\"ursey
symmetries of Dirac matter, and false for every spacetime symmetry, the six-dimensional one included: a mass changes which little group a state has
and removes the duality generator from the stabiliser of a field.

\emph{$c^2$ is the symmetric part} is an interpretive entry, and it needs correction. The constant $c$ is the Schur-forced conversion $C=c\,\id$
between the triples (Theorem~3.2 of \cite{ES17}), and a constant generates nothing. In the interior it is carried by $\kappa G$, $\kappa=c/\ell$, where
$c^2$ appears as $(\kappa G)^2=-c^2/\ell^2$, the squared binding frequency; and $G$ generates the broken duality. But $G$ is not an infinitesimal isometry
of $\eta$ (\cref{d2:prop:duality}): its duality is a symmetry of the positive form $u$, while the interior Lagrangian is built on $\eta$. These are the
duality and Lorentz classes of Theorem~17.11 of \cite{ES17}, which exclude each other, so the symmetric part and the Lorentz-invariant dynamics cannot
be one object. In the Dirac reading the symmetric part is the massless kinetic algebra $\uu1_\Omega\oplus\su2$, which the mass cuts to $\uu1_J$; for
$\mathrm{Spin}(3,3)$ the roles are reversed (\cref{d2:prop:Tinv}). Which term is symmetric depends on which symmetry is meant, and the hypothesis has to
say which. Nor is the undivided energy the interior's conserved energy: $u$ is frame-dependent and not conserved by the binding flow, the Noether
energy has duality charge two and lives on the overlap of the two senses, and the only listed quantity invariant under both Lorentz and duality is
$E_0=\half|q|$, the infimum of $u$ over Lorentz frames (\cref{d2:prop:rest,d2:prop:interior}). If an undivided energy whose differentiation is the mass
is wanted, the arena's candidate is $u$ with its split $u^2=E_0^2+|\mathbf S|^2$, the inherited field form of $E^2=p^2c^2+m^2c^4$; the differentiation it
describes is the choice of a rest frame, not a breaking.

Finally the $b$-field. At one constituent no broken symmetry can carry a Yang--Mills field: duality has no spin lift, $U(1)_\Omega$ is anomalous, and the
Pauli--G\"ursey $\su2$ acts on a single doublet and mixes a field with its conjugate (Proposition~5.5 of \cite{ES17}, \cite{Witten82}). The one gaugeable
generator there, $J$, is abelian and unbroken by the Dirac mass ($[T,J]=0$); it acquires a St\"uckelberg mass only from a vacuum projector that fails to
commute with it (Theorem~26.1 of \cite{ES17}). In the paper's vacuum analysis broken non-abelian generators are first gauged at three constituents,
where the nine Goldstone directions of $\Phi_0$ are eaten by $Z'$, $W_R^\pm$ and six leptoquarks (\cref{d2:prop:goldstone}; Theorems~21.1 and~22.2 of
\cite{ES17}). There the gauge quanta of Yang and Mills \cite{YangMills54} acquire their masses from the breaking, by the inherited mechanism of
\cite{EnglertBrout64,Higgs64}, and the fermion masses come from the Yukawa coupling to the same condensate. In the arena the $b$-particle receives mass
from the differentiation; it does not generate it. That is where the $b$-field belongs, and it reverses the order the hypothesis proposes.

\begin{table}[t]
\centering
\small
\begin{tabular}{@{}lll@{}}
\toprule
statement & class & checks\\
\midrule
\cref{d2:prop:duality} (i), (iii), (iv) / (ii) & inherited / fixed & d2-1 / d2-2\\
\cref{d2:prop:rest} & inherited & d2-3\\
\cref{d2:prop:interior} & fixed & d2-4\\
\cref{d2:thm:little} & inherited & d2-5a, d2-5, d2-6, d2-7\\
\cref{d2:prop:Tinv} & fixed & d2-8, d2-8b\\
\cref{d2:prop:massless} (i), (ii) / (iii), (iv) & fixed / inherited & d2-9 / d2-10\\
\cref{d2:prop:restflow} & inherited & d2-11\\
\cref{d2:prop:senses} (i) / (ii), (iii) & inherited / fixed & d2-12 / d2-13, d2-14\\
\cref{d2:prop:fock} & inherited & d2-15a, d2-15b, d2-16, d2-17\\
\cref{d2:prop:framevac} / \cref{d2:cor:noground} & fixed / obstructed & d2-18 / d2-19\\
\cref{d2:prop:goldstone} & fixed & d2-20, d2-21a, d2-21, d2-22\\
electroweak axion-like count (Theorem~22.3 of \cite{ES17}) & fixed & cited, not re-verified\\
$(\kappa G)^2=-c^2/\ell^2$; readings of \cref{d2:sec:verdict} & fixed; interpretive & d2-23; ---\\
\bottomrule
\end{tabular}
\caption{Claims of this section, their class, and the labels of the checks printed by \texttt{d2\_symmetry\_breaking.py}; check d2-0 verifies the
realisation.}
\label{d2:tab:claims}
\end{table}
